% Insert as A.65.  A.66 supplies the arithmetic construction used below.
\begin{statement}{A.65}{Proved}{Ordinary infinite-support ghosts at every fixed nilpotent depth}
For every integer $R\geq1$ and every positive integer $N$, there is an
ordinary completed unit-equation pair $(A,B)$ with infinite valuation
support, together with an actual section of
\[
 E_R=\pi_1(U)/\Gamma_{R+1}\longrightarrow G_\Q,
 \tag{A.65.1}
\]
whose abelianization is $(A,B)$.  This section can be chosen to agree with
the section of the rational point $3$ in the finite quotient obtained by
reducing the truncated integral Magnus expansion through degree $R$
modulo $N$.  Its abelianization is neither a rational-point image nor a
cuspidal image.

The pair and the section constructed here may depend on $R$ and $N$.
The assertion supplies actual full integral sections at each fixed
nilpotency class; it does not assert a single pair or a compatible
sequence of sections as $R$ varies.
\end{statement}
\begin{proof}
It suffices to prove the assertion for $R\geq3$: the resulting sections
project to classes one and two.  We retain a fixed rational tangential
splitting and a fixed representative $s_0$ of the rational point $3$.
Relative geometric cocycles will be taken with respect to $s_0$.

\emph{Integral Hall coordinates and the exact finite congruence.}
Write $N_r=\Pi/\Gamma_{r+1}$.  Choose integral Hall commutators of weights
at most $R$, and let $h_R$ be their number.  One can, equivalently, use
the standard bracketings of Lyndon words.  At each weight their
lowest-degree Magnus polynomials have a unit triangular coefficient
matrix on the corresponding Lyndon words.  The elementary induction
behind this assertion brackets the two standard factors of a Lyndon
word: their concatenation occurs with coefficient one, and every other
word is larger in lexicographic order.  Antisymmetry and Jacobi collect
an arbitrary commutator into these bracketings.  Thus these polynomials
give an integral basis and an integral coordinate-extraction map at
each weight.

Collecting successively in the central quotients gives unique continuous
Hall coordinates in $N_r$, with coefficients in $\widehat\Z$.  In
particular
\[
 L_r=\Gamma_r/\Gamma_{r+1}
       \simeq\widehat\Z(r)^{\,b_r},\qquad
 b_r=\#\{\text{Hall commutators of weight }r\}.
 \tag{A.65.2}
\]
The generating assertion follows by commutator collection; independence
follows from the unit triangular Magnus matrices just described.  This
also proves the corresponding assertions after projection to any
$\Z_\ell$ component.

Let $I$ be the positive-degree ideal in the noncommutative polynomial
ring, truncated by $I^{R+1}$.  The map $x\mapsto1+X$, $y\mapsto1+Y$
defines the actual truncated Magnus homomorphism.  Define
\[
 K_r(m)=
 \ker\!\left(N_r\longrightarrow
      \bigl(1+I\bigr)
       \bigl((\Z/m)\langle X,Y\rangle/I^{r+1}\bigr)\right).
 \tag{A.65.3}
\]
Integral substitution of the two generators shows that this kernel is
preserved by every continuous automorphism of $N_r$.  Hence it defines
a finite geometric quotient.  More explicitly, write
$E_r=N_r\rtimes_{\rho_0}G_\Q$ using $s_0$ and let
$\bar\rho_0$ be its action on $N_r/K_r(m)$.
The actual finite arithmetic quotient used in the statement is the map
$(z,g)\mapsto(z\bmod K_r(m),\bar\rho_0(g))$ to
$(N_r/K_r(m))\rtimes\operatorname{im}(\bar\rho_0)$.
Agreement with $s_0$ modulo $m$ will mean the pointwise assertion
$s(g)s_0(g)^{-1}\in K_r(m)$, which gives equality under this map.
No identification of
$\ker(N_r/K_r(m)\to N_{r-1}/K_{r-1}(m))$
with $L_r/m$ is made.

There is nevertheless a controlled actual set lift.  Let
$\sigma_r:N_{r-1}\to N_r$ keep the lower Hall coordinates and set the
weight-$r$ coordinates to zero.  Put
\[
       D=(R!)^{h_R+1},\qquad T=2(R!)^R.
 \tag{A.65.4}
\]
Then
\[
             z\in K_{r-1}(Dm)
        \quad\Longrightarrow\quad \sigma_r(z)\in K_r(m).
 \tag{A.65.5}
\]
Here and below the same bounds are valid in each $\Z_\ell$ component.
For completeness, extract the Hall coordinates from the successive
lowest-degree Magnus coefficients and remove each extracted Hall
power.  If an extracted exponent $a$ is divisible by an integral ideal
$J$, the nonconstant coefficients of that Hall power through degree $R$
belong to $J/(R!)$, since
$\binom ak=(a/k)\binom{a-1}{k-1}$ for $k\leq R$.
The coordinate-extraction matrices are integral, so each removal loses
at most this one factor $R!$.  There are at most $h_R$ removals.
Starting with $Dm$ leaves all coefficients of the resulting canonical
lift divisible by $m$.  All the intermediate ideals are integral.
This proves (A.65.5), including the small coefficient primes.

\emph{A finite annihilator of the unramified obstruction.}
For $r\geq2$ set
\[
 B_r=2\,|K_{2r-2}(\Z)|,\qquad
 D_r^{\rm old}=(2^{r-1}-1)(3^{r-1}-1).
 \tag{A.65.6}
\]
The even positive $K$-group here is finite.  The precise comparison
needed below is
\[
 B_r\,H^2_{\rm cts}
       (G_{\Q,\{\ell,\infty\}},\Z_\ell(r))=0
                  \quad\hbox{for every prime }\ell.
 \tag{A.65.7}
\]
For odd $\ell$, the group is the $\ell$-primary part of
$K_{2r-2}(\Z)$ by \cite[VI.8.2]{Weibel2013}.
Localization at $\ell$ does not change that primary part: the adjacent
finite-field $K$-groups have no $\ell$-torsion.
For $\ell=2$, apply \cite[VI.9.11]{Weibel2013} to $\Z[1/2]$.
Here $r_1=1$ and the integer $\rho$ in that theorem is zero.
For $r\not\equiv0\pmod4$ the indicated cohomology group is the
corresponding two-primary $K$-group.  For $r\equiv0\pmod4$, its kernel
under real restriction is that $K$-group and its remaining image is a
subgroup of $\Z/2$.  The extra factor two in (A.65.6) therefore gives
(A.65.7) in every case.  Finiteness follows from the finite-generation
and rank results recorded in \cite[IV.1.18, IV.6.9, VI.8.1]{Weibel2013}.
No uncomputed even $K$-group is being declared zero.

We also use the actual continuous localization map, in each weight:
\[
 \begin{split}
 H^2_{\rm cts}(G_{\Q,\{\ell,\infty\}},\Z_\ell(r))
   &\longrightarrow H^2_{\rm cts}(G_\Q,\Z_\ell(r))\\
   &\longrightarrow
 R_\ell/(q^{r-1}-1)R_\ell,\qquad
 R_\ell=\{(z_q)_{q\ne\ell}:z_q\in\Z_\ell,\ z_q\to0\}.
 \end{split}
 \tag{A.65.8}
\]
It is exact at the middle term.  This is the strict localization
complex of A.27 with the coefficient twist changed to $r$:
its two local terms are $\Z_\ell(r-1)$ with differential
$q^{r-1}-1$.  The coefficient-compatible tame averaging and
normalization maps used there give the same inverse-limit complex.
In particular the image in (A.65.8) uses the literal module $R_\ell$,
not its quotient by a closure.

\emph{Uniform control of the actual fresh inertia.}
At $q\ne\ell$, the tangential action on every $N_r^{(\ell)}$ is
unramified.  The class-two assertion is the actual Ihara formula.
Inductively, the kernel of the automorphism map from class $i+1$
to class $i$ is
$\operatorname{Hom}(N_1,L_{i+1})$.
Frobenius conjugation on it is $q^i$, whereas the tame relation acts
on the inertia generator by $q$.  Thus an inertia image $\delta$
satisfies $(q^i-q)\delta=0$ for $i\geq2$ and is zero in this
torsion-free module.  Wild inertia and the other tame primary factors
have zero image as well.  This proves the induction.

The geometric cocycle of $s_0$ is unramified at every $q>3$:
its abelian Kummer coordinates are $3,-2$, which have zero valuations
there, and the same tame relation successively kills every higher
central inertia coordinate.  Consequently the action relative to
$s_0$ is also unramified at these primes.

Consider an actual section through class $i$, at a prime $q>3$ with
ordinary abelian valuation pattern $(E,0)$, $(0,E)$, or $(-E,-E)$.
Write its tame inertia in Hall coordinates, and put
$V_2=y^n x^m C^t$.  The ordinary local unit equation gives the exact
class-two identities used in A.27 and A.58:
\[
       m,n\in E\Z_\ell,\qquad
       t\in(E/2)\Z_\ell .
 \tag{A.65.9}
\]
For the two positive patterns $t$ is the valuation times an integral
unit parameter, up to sign; for the negative pattern it is
$\binom m2-mu$ with $u\in\Z_\ell$.
Thus the canonical Hall lift of $V_2$ through degree $R$ is one modulo
$E/(2R!)$.

Suppose the canonical lift of the inertia through weight $i-1$ is
$\bar V$, and the actual weight-$i$ coordinates form $Z_i\in L_i$.
The exact tame relation gives
\[
 U\varphi_q(\bar V)U^{-1}\bar V^{-q}
                    =Z_i^{-q(q^{i-1}-1)}
                         \quad\hbox{in class }i.
 \tag{A.65.10}
\]
Integral Frobenius substitution, conjugation and powers preserve every
integral Magnus coefficient ideal.  The left side is central, so the
primitive weight-$i$ extraction matrix introduces no denominator.
Division by $q^{i-1}-1$ bounds its actual central coordinates.
Extending the Hall powers of those coordinates through degree $R$
loses at most one further factor $R!$.  Induction therefore bounds the
canonical lift of the inertia through class $i$ by the ideal
\[
 \frac{E}{\,2(R!)^{i-1}\prod_{k=2}^{i-1}(q^k-1)}\Z_\ell
                         \quad(2\leq i\leq R).
 \tag{A.65.11}
\]
All these ideals will be integral for the exponents chosen below.

In particular, for the normalized defect $c$ of a canonical lift from
class $r-1$ to class $r$, its actual strict tame residue satisfies
\[
 \frac{\psi_q(c)}{q^{r-1}-1}
       \ \in\
 \frac{E}{\,T\prod_{k=2}^{r-1}(q^k-1)}
          \Z_\ell^{\,b_r}.
 \tag{A.65.12}
\]
For $r\geq3$, this follows from (A.65.11) and the exact equality
$\psi_q(c)=q^{-1}\mathcal R_q$, where $\mathcal R_q$ is the lifted
tame relator, as in A.61.  For $r=2$, the original-cup residue in A.27,
or the same class-two relation with the ordinary parameter in
(A.65.9), gives $\psi_q(c)/(q-1)\in(E/2)\Z_\ell$.
This is stronger than (A.65.12), with the product interpreted as empty.
Changing from the tangential splitting to $s_0$ does not alter this
calculation: at $q>3$ its inertia is one, and its Frobenius is absorbed
in the integral element $U$ and the actual relative action.

When both valuations are zero, inertia is one in every nilpotent
quotient, by the torsion-free tame induction.  For the canonical lift,
the strict residue is then literally zero.

\emph{The exact central correction preserving a finite congruence.}
Set $M_R=N$, and define the positive integers $M_{r-1}$ backwards by
\[
 M_{r-1}=D\,B_r M_rD_r^{\rm old}\qquad(2\leq r\leq R).
 \tag{A.65.13}
\]
Choose a positive integer $C$ divisible by $24$, by $M_1$, and by
$T B_rM_r$ for every $r=2,\ldots,R$.
The construction A.66, with $d=R-1$, provides an ordinary pair with
fixed valuations $(0,1)$ at two and $(1,0)$ at three, and, at each
fresh prime $q>3$, one of the patterns with exponent
\[
 E_q=C e_q\prod_{k=2}^{R-1}(q^k-1),\qquad e_q\longrightarrow0
                                  \text{ in }\widehat\Z.
 \tag{A.65.14}
\]
It also gives $A/3$ and $B/(-2)$ as $M_1$th powers in the global power
completion.  Hence we can choose the abelian relative cocycle
pointwise divisible by $M_1$.  This is an actual class-one section
agreeing with $s_0$ modulo $M_1$.

Assume the section through class $r-1$ has been constructed with this
agreement modulo $M_{r-1}$.
Use $\sigma_r$ on its relative geometric cocycle.
By (A.65.5), its set lift agrees pointwise with $s_0$ through degree
$r$ modulo
\[
                   L=B_rM_rD_r^{\rm old}.
 \tag{A.65.15}
\]
Its defect is central.  Integral substitution and primitive central
coordinate extraction show that the normalized defect cocycle $c$
takes values in $L\widehat\Z(r)^{b_r}$.
Thus $c/(B_rM_r)$ is an actual continuous integral two-cocycle.

Its strict residues divided by $q^{r-1}-1$ are integral at two and
three, because (A.65.15) contains their two denominators.
Averaging and normalizing local cochains preserve this coefficient
divisibility.  At every fresh prime, (A.65.12)--(A.65.14) put the
quotients in $e_q\Z_\ell^{b_r}$, by the choice of $C$.
At all other $q>3$ they vanish literally.
Consequently these quotient families belong to $R_\ell^{b_r}$.
Exactness of (A.65.8) places every component of
\[
                  \beta=[c/(B_rM_r)]
 \tag{A.65.16}
\]
in the image of unramified cohomology.
By (A.65.7), $B_r\beta=0$, including the real two-primary contribution.
In particular $[c]=B_rM_r\beta=0$.
Choose a continuous integral one-cochain $h$ with $\delta h=-c$.
Adding $h$ gives an actual class-$r$ section, initially without a
congruence assertion on its highest coordinate.

The needed correction follows from an exact coefficient diagram.
Since $c$ is divisible by $B_rM_r$, the reduction of $h$ is a
one-cocycle with coefficients $\Z/(B_rM_r)(r)^{b_r}$, and
\[
 \partial_{B_rM_r}[\bar h]=-\beta,\qquad
 \partial_{M_r}[\bar h\bmod M_r]
                =B_r\partial_{B_rM_r}[\bar h]
                =-B_r\beta=0.
 \tag{A.65.17}
\]
The second equality is the cochain identity
$\delta h/M_r=B_r(\delta h/(B_rM_r))$.
Thus $\bar h\bmod M_r$ lifts to an integral continuous one-cocycle
$z$ with values in $L_r$.  This is the long exact sequence for
$0\to L_r\xrightarrow{M_r}L_r\to L_r/M_r\to0$.
If the reductions initially agree only as cohomology classes,
their difference is a finite coboundary; lift its zero-cochain to
$L_r$ and subtract its integral coboundary from $z$.
We then have the pointwise equality
$z\equiv h\pmod{M_r}$.
Replacing $h$ by $h-z$ preserves $\delta h=-c$ and makes every highest
coordinate divisible by $M_r$.  Together with (A.65.15) this yields
the required actual agreement with $s_0$ in (A.65.3).
This completes the induction.

All coefficient-prime constructions combine: continuous cochains with
a product of compact coefficient modules are the product of the
individual cochain complexes, and multiplication by any nonzero
integer is a homeomorphism onto its image.
The resulting section is therefore a section of the full integral
quotient (A.65.1).

Finally A.66 gives infinitely many distinct nonzero valuation
coordinates in each of the three support patterns.
No rational number has that valuation support.
A cusp at zero has trivial $B$ coordinate, a cusp at one has trivial
$A$ coordinate, and a cusp at infinity has equal two valuation vectors;
the fixed valuations at two and three exclude all these images.
Thus the constructed section has the asserted non-point,
noncuspidal abelianization.
\end{proof}

